\begin{hblemma}{thm:p12-har-fw-normalization}{proved=annex:p12.har.fw-normalization}{The exact LF host has typed beta-normal representatives}
Every derivable judgment over a well-formed signature/context in this
host has a beta-normal subject of the same type. The identity sort/rule
inclusion in exact CC gives strong normalization; the source rule
construction gives LF subject reduction. Normal-form generation exposes
the head constant/variable and its complete typed application spine.
Opaque atomic families \(Tq\), \(\mathrm{eq}\,uv\),
\(\mathrm{prop}\), \(A\), and \(D\) cannot become dependent
products by beta conversion. This is not an assumption that every
normal term is beta-convertible to an eta-long canonical term.
\end{hblemma}

\begin{hbdefinition}{def:p12-har-fw-contexts}{Admissible contexts and pointwise reification}
For PROP, the context after \(\Sigma_{\mathrm{PROP}}\) consists
of finitely many distinct atom names \(a_i:\mathrm{prop}\) and
proof variables \(h_j:T(q_j)\). For PRED, after the full signature
(including inherited \(T\)), it consists of object variables
\(x_i:A\) and proof variables \(h_j:T(q_j)\).
The \(q_j\) are well-typed formula names, normalized when decoding.
For the untyped-lambda converse, the context after
\(\Sigma_{\mathrm{lambda}}\) consists only of free object
variables \(x_i:D\); an extra equality-proof assumption would assert
an extra equation, not pure beta equality.

Normalize before reification. PROP formula names are atoms or
\(\mathrm{imp}\,p q\). PRED object names are variables, \(c\),
or \(ft\); formula names are \(Rst\), \(\mathrm{imp}\,p q\),
or \(\mathrm{forall}\,F\). Decode the first two formula cases
structurally. For the third, decode
\(\mathrm{forall}\,F\) as
\(\forall z:A.\mathcal D(\mathrm{nf}_\beta(Fz))\), with \(z\)
fresh. A normal \(F:A\to\mathrm{prop}\) is either
\(\lambda z:A.q\) or \(Rs\); the latter case directly yields
\(\forall z:A.R(\mathcal D(s),z)\).

For normal untyped-lambda names, decode \(x:D\) to the source
variable, \(\mathrm{app}\,uv\) to the source application, and
\(\mathrm{abs}\,F\) to
\(\lambda z.\mathcal D(\mathrm{nf}_\beta(Fz))\).
A normal \(F:D\to D\) is either \(\lambda z:D.u\) or
\(\mathrm{app}\,u\); the latter directly yields
\(\lambda z.\mathcal D(u)z\). These pointwise clauses handle
eta-short LF functions without adding an object-language eta equation.
For example \(\mathrm{abs}(\mathrm{app}\,u)\) decodes to
\(\lambda z.\mathcal D(u)z\), not to \(\mathcal D(u)\).
Freshness is relative to the entire current context and the arguments.
\end{hbdefinition}

\begin{hblemma}{thm:p12-har-fw-reification}{proved=annex:p12.har.fw-reification}{The reifiers are defined and commute with open substitution}
In the stated contexts, the normal-name grammars above are exhaustive,
independent of proof variables, and their reifiers terminate. Reification
of the source encoding recovers each source expression up to bound-name
renaming. Reification is invariant under LF beta conversion and commutes
with substitutions of object names for object variables, including
capture-avoiding substitution under \(\mathrm{forall}\) and
\(\mathrm{abs}\). In PRED these are arbitrary first-order term
substitutions, not assignments only to closed constants; in the lambda
case they are arbitrary finite open untyped-term substitutions.
\end{hblemma}

\begin{hbtheorem}{thm:p12-har-fw-faithfulness}{proved=annex:p12.har.fw-faithfulness}{Derivability reflection for all three source framework signatures}
At exactly the host/signature/context discipline above:
\begin{enumerate}
 \item If \(\Sigma_{\mathrm{PROP}},\vec a,\vec h
 \vdash p:T([\varphi])\), then the decoded assumptions
 \(\mathcal D(q_j)\) derive \(\varphi\) in minimal implication
 natural deduction. With no \(\vec h\), \(\varphi\) is a theorem.
 \item If \(\Sigma_{\mathrm{PRED}},\vec x,\vec h
 \vdash p:T([\varphi])\), then those decoded assumptions derive
 \(\varphi\) in the specified one-sorted minimal
 \(\{\supset,\forall\}\)-logic with \(c,f,R\).
 The constructed universal-introduction rule obeys its eigenvariable
 condition, and elimination uses the actual decoded open argument.
 \item If \(\Sigma_{\mathrm{lambda}},\vec x:D
 \vdash p:\mathrm{eq}[P][Q]\), then \(P=_\beta Q\)
 in the untyped lambda calculus, with its actual open free variables.
\end{enumerate}
The proof constructs a source derivation from every normal proof head,
including function-valued proof arguments to introduction/xi constants.
Combined with \cref{thm:p12-har-framework-forward}, these give both
existence directions. They do not assert a bijection on all LF terms or
faithfulness after adding arbitrary signature axioms
\cite[pp.~1161--1163]{BarendregtGeuvers2001}.
\end{hbtheorem}

\begin{hbcorollary}{thm:p12-har-fw-no-eta}{proved=annex:p12.har.fw-no-eta}{The lambda signature has not acquired object eta}
In the pure lambda signature with \(u:D\), no term inhabits
\(\mathrm{eq}\,u\,[\lambda x.ux]\), for \(x\) fresh.
Pointwise reification of LF function arguments is not the object
equation \(u=\lambda x.ux\).
\end{hbcorollary}


